\documentclass[12pt]{article}
\usepackage[T1]{fontenc}
\usepackage[utf8]{inputenc}
\usepackage[brazil]{babel}
\usepackage[margin=1in]{geometry}
\usepackage{ragged2e}
\usepackage[normalem]{ulem}
\setlength{\parindent}{0pt}
\setlength{\parskip}{0.8\baselineskip}
\begin{document}
\justifying
\noindent Verifica-se que a questão objeto da controvérsia jurídica suscitada no recurso manejado pela parte recorrente esteve afetada no representativo da controvérsia - \textbf{Tema n. }\textbf{243} da TNU - (PU no processo n. PEDILEF 001423889-2015.4.01.3700/MA),  afetado em 06/11/2019 foi julgado (trânsito em julgado em PUIL/STJ nº 2712/MA), por meio do qual foi elucidada a seguinte questão:

\noindent \textit{"Saber se a demora excessiva na fila de atendimento em instituição financeira enseja indenização por dano moral"}

\noindent Restou firmada a seguinte tese:

\noindent \textit{"I) a espera em fila de banco por tempo superior ao previsto na legislação local não configura, por si só, dano moral in re ipsa; II) é cabível indenização por danos morais fundada na espera em fila de banco quando a demora for excessivamente longa ou quando estiver associada a outros constrangimentos capazes de abalar a honra do autor ou causar-lhe situação de dor, sofrimento ou humilhação." (Tema 243 TNU).}

\noindent Com efeito, depreende-se da leitura do voto condutor do acórdão que o mesmo se encontra em conformidade com o decidido no  tema pela TNU.

\noindent A proposito, consta do acordao hostilizado: (...) Financiamento Privado do Ensino Superior e/ou Pesquisa, FINANCIAMENTO, DIREITO À EDUCAÇÃO, Indenização por dano moral, Responsabilidade civil, DIREITO CIVIL, Nulidade, Atos Processuais, DIREITO PROCESSUAL CIVIL E DO TRABALHO

\noindent Nessas condições, o conteúdo do Acórdão recorrido é consentâneo com o entendimento da TNU, firmado em julgamento de recurso representativo de controvérsia, o que atrai a incidência da regra prevista pelo art. 14, III, \textit{b,} da Resolução CJF n. 586/2019:

\noindent \textit{\textbf{Art. 14.}}\textit{ Decorrido o prazo para contrarrazões, os autos serão conclusos ao magistrado responsável pelo exame preliminar de admissibilidade, que deverá, de forma sucessiva: (...) }\textit{\textbf{III }}\textit{- negar seguimento a pedido de uniformização de interpretação de lei federal interposto contra acórdão que esteja em conformidade com entendimento consolidado: }\textit{\textbf{b)}}\textit{ em }\uline{\textit{\textbf{recurso representativo de controvérsia pela Turma Nacional de Uniformização}}}\textit{ ou em pedido de uniformização de interpretação de lei dirigido ao Superior Tribunal de Justiça;(...).}

\noindent Posto isso, com arrimo no \uline{\textbf{art. 14, III, }}\uline{\textit{\textbf{b}}}\textit{,} da Resolução CJF n. 586/2019 c/c art. 1.030, I, do CPC, \textbf{NEGO SEGUIMENTO }\textbf{a}o\textbf{ PU}\textbf{.}

\noindent Intimem-se. Aguarde-se o decurso do prazo recursal. Após, certifique-se o trânsito em julgado e devolva-se ao juízo de origem.
\end{document}
